\documentclass[11pt,aspectratio=169]{beamer}
\usetheme{Madrid}
\usecolortheme{default}
\usepackage{graphicx}
\usepackage{booktabs}
\usepackage{hyperref}
\usepackage{xcolor}

% Title Slide Information
\title[Code-Switched RAG for Nepali PDFs]{Overcoming Font-Encoding Challenges in Nepali-English Code-Switched RAG Systems}
\subtitle{A Local LangChain Retrieval Evaluation}
\author[Adhikari Pankaj]{ADHIKARI PANKAJ \\ \vspace{0.2cm} \small Student ID: 2120256037 \\ Masters Student}
\institute[Nankai University]{Nankai University}
\date{\today}

\begin{document}

% Slide 1: Title Page
\begin{frame}
    \titlepage
\end{frame}

% Slide 2: Background
\section{Background}
\begin{frame}{Background}
    \begin{itemize}
        \item \textbf{Data Accessibility in Nepal:} Crucial municipal documents (e.g., tax schedules, business registrations) are heavily siloed in legacy PDF formats.
        \item \textbf{The RAG Promise:} Retrieval-Augmented Generation (RAG) can help citizens query these documents naturally using code-switched language (Nepali + English).
        \item \textbf{The Hidden Barrier:} Many legacy South Asian PDFs utilize custom font-encodings (like Preeti) instead of standard Unicode. 
        \item \textbf{The Resulting Mess:} When extracted, standard text appears as gibberish (e.g., \textit{"sf7df8f\}F dxfgu/kflnsf"} instead of actual Nepali).
    \end{itemize}
\end{frame}

% Slide 3: Motivations
\section{Motivations}
\begin{frame}{Motivations}
    \begin{itemize}
        \item \textbf{Catastrophic Baseline Failure:} State-of-the-art dense embedding models (like BGE-M3) return a \textbf{0.000 Recall/MRR} on this raw data.
        \item \textbf{Misleading Evaluations:} Standard AI evaluations might falsely conclude the AI model is bad, missing the root cause: font-encoding corruption.
        \item \textbf{Practical Need:} We need a robust, local retrieval fallback strategy. A 0.000 result is not a dead end; it is a data pipeline revelation that demands metadata solutions.
    \end{itemize}
\end{frame}

% Slide 4: Aim & Objective
\section{Aim \& Objective}
\begin{frame}{Aim \& Objective}
    \begin{block}{Primary Aim}
        To design and evaluate a local LangChain-based hybrid retrieval system capable of accurately retrieving information from corrupted Nepali PDFs using code-switched queries.
    \end{block}
    \vspace{0.3cm}
    \textbf{Specific Objectives:}
    \begin{enumerate}
        \item Measure the exact failure rates of embedding models on font-encoded PDFs.
        \item Build a project-owned chunk database augmenting corrupted text with normalized English and Unicode aliases.
        \item Evaluate chunk-level retrieval using a local LangChain hybrid retriever.
    \end{enumerate}
\end{frame}

% Slide 5: Literature Review
\section{Literature Review}
\begin{frame}{Literature Review}
    \begin{itemize}
        \item \textbf{Low-Resource RAG:} Models trained primarily on English struggle with cross-lingual and code-switched tasks due to tokenizer misalignment (Ahuja et al., 2023).
        \item \textbf{Document Digitization in South Asia:} Legacy font encoding remains a massive bottleneck for NLP in languages like Nepali, Hindi, and Bengali.
        \item \textbf{Hybrid Retrieval Systems:} Combining sparse retrieval (TF-IDF/BM25) with dense embeddings and metadata augmentation is proven to significantly improve Out-Of-Vocabulary (OOV) query resolution (Gao et al., 2023).
    \end{itemize}
\end{frame}

% Slide 6: Methods
\section{Methods}
\begin{frame}{Methods: The Local LangChain RAG System}
    \begin{itemize}
        \item \textbf{Dataset:} 446 extracted chunks from 3 municipal PDFs (Lalitpur, Bhaktapur, Kathmandu).
        \item \textbf{Query Bank:} 11 real-world code-switched Nepali-English queries.
        \item \textbf{Pipeline Construction:} 
        \begin{itemize}
            \item Built a local LangChain Core retrieval chain.
            \item Replaced standard dense embeddings with a persistent local project-owned database.
            \item \textbf{Metadata Augmentation:} Injected normalized aliases (Romanized Nepali + English keywords) into the chunk metadata to bridge the OCR gap.
        \end{itemize}
    \end{itemize}
\end{frame}

% Slide 7: Evaluation Metrics Explained
\section{Understanding Metrics}
\begin{frame}{Understanding the Evaluation Metrics}
    \begin{block}{What is Recall@5?}
        \textbf{Recall@5} measures the probability that the exact correct document appears in the \textbf{top 5 search results}.
        \begin{itemize}
            \item \textit{Why it matters:} In RAG, the LLM reads the top results. If the factual answer is not in the top 5, the LLM will hallucinate.
        \end{itemize}
    \end{block}
    
    \vspace{0.3cm}
    \begin{block}{What is Mean Reciprocal Rank (MRR)?}
        \textbf{MRR} evaluates \textbf{where} the correct answer ranks.
        \begin{itemize}
            \item If the correct answer is rank 1, score = 1.0. Rank 2 = 0.5. Rank 3 = 0.33.
            \item \textit{Why it matters:} Higher MRR means the user (or LLM) doesn't have to read through irrelevant documents to find the truth.
        \end{itemize}
    \end{block}
\end{frame}

% Slide 8: Results
\section{Results}
\begin{frame}{Results \& Findings}
    \begin{table}[]
        \centering
        \begin{tabular}{@{}lcc@{}}
        \toprule
        \textbf{System Configuration} & \textbf{Recall@3} & \textbf{MRR} \\ \midrule
        Dense Baseline (BGE-M3 on raw text) & 0.000 & 0.000 \\ 
        Raw Sparse (TF-IDF on raw text) & 0.180 & 0.150 \\ 
        \textbf{Local RAG (LangChain + Aliasing)} & \textbf{0.917} & \textbf{0.917} \\ \bottomrule
        \end{tabular}
    \end{table}
    \begin{itemize}
        \item \textbf{Observation:} Pure semantic matching completely fails on font-corrupted text.
        \item \textbf{Impact:} By augmenting chunks with normalized semantic aliases, the local LangChain retriever jumps from 0\% to \textbf{91.7\% accuracy}, achieving a highly usable system.
    \end{itemize}
\end{frame}

% Slide 9: LLM vs RAG Example
\section{Direct LLM vs Local RAG}
\begin{frame}{Direct LLM vs. Our Local LangChain System}
    \begin{table}[]
        \centering
        \footnotesize
        \begin{tabular}{@{}p{3.8cm} p{4.5cm} p{4.5cm}@{}}
        \toprule
        \textbf{Citizen Query} & \textbf{Direct LLM (ChatGPT/Claude)} & \textbf{Our Local LangChain RAG} \\ \midrule
        
        \textit{"Kathmandu public parking ma two-wheeler fee kati?"} & 
        \textcolor{red}{\textbf{Hallucination:}} ``Typically, it is Rs 20/hr... check the website.'' & 
        \textcolor{green}{\textbf{Accurate (from PDF):}} ``As per KTM\_FIN\_002, the fee is Rs 15 for the first half-hour.'' \\ \midrule
        
        \textit{"Property tax slab ma 1-2 crore rate kati? exact percent chahincha."} & 
        \textcolor{red}{\textbf{Missing Knowledge:}} ``I don't have access to the exact 2024 tax slabs.'' & 
        \textcolor{green}{\textbf{Accurate (from PDF):}} ``The property tax for the 1 to 2 crore slab is 0.015\% (Rs 1500).'' \\ \midrule
        
        \textit{"Bhaktapur naksa pass application documents list?"} & 
        \textcolor{red}{\textbf{Generic Guess:}} ``You generally need citizenship, land document, and blueprint.'' & 
        \textcolor{green}{\textbf{Accurate (from PDF):}} Retrieves the exact 11-item checklist from Section 3. \\ \bottomrule
        \end{tabular}
    \end{table}
\end{frame}

% Slide 10: Outline (For the Paper)
\section{Outline}
\begin{frame}{Outline (For the Paper)}
    \begin{enumerate}
        \item \textbf{Introduction:} The digitization gap in Nepali municipalities.
        \item \textbf{Related Work:} Code-switching, Hybrid Retrieval, Legacy Encodings.
        \item \textbf{Methodology:} 
        \begin{itemize}
            \item Dataset creation and font-encoding analysis.
            \item The Alias-Augmented Local LangChain Retriever.
        \end{itemize}
        \item \textbf{Experiments \& Results:} Baseline vs. Augmented Performance metrics.
        \item \textbf{Discussion:} RAG evaluations vs Direct LLM.
        \item \textbf{Conclusion \& Future Work:} Integrating vision-language models (VLMs) bypassing OCR.
    \end{enumerate}
\end{frame}

% Slide 11: Conclusion
\section{Conclusion}
\begin{frame}{Conclusion}
    \begin{center}
        \Large \textbf{Thank You!}
        \vspace{0.5cm}
        
        \normalsize
        The initial 0.000 baseline was not a dead end, but a revelation about the data pipeline. By applying smart metadata augmentation and a local LangChain retrieval chain, we transformed inaccessible municipal PDFs into a queryable, high-performing RAG corpus.
        
        \vspace{0.5cm}
        \textit{Questions?}
    \end{center}
\end{frame}

\end{document}
